\documentclass[slideColor]{prosper}
\usepackage{amsmath,amsfonts,amsthm}
%\usepackage{epsfig}
%\include{Qcircuit}

%\usepackage{algorithm} 
%\usepackage{algorithmic}

\newcommand{\E}{\mathbf{E}}
\renewcommand{\R}{\mathbb{R}}
\newcommand{\Z}{\mathbb{Z}}

\newcommand{\cA}{\mathcal A}
\newcommand{\cB}{\mathcal B}
\newcommand{\cH}{\mathcal H}
\newcommand{\C}{\mathbb C}
\newcommand{\TVD}{\mathrm{TVD}}

\renewcommand{\>}{\rangle}
\newcommand{\<}{\langle}   


%%%%%%%%%%%%%%%%%%%%%%%%%%%%%%%%%%%%%%%%%%%%%%%%%%%%%%%%%%%%

\title{Quantum Algorithm for Computing the Period Lattice of an Infrastructure}
\author{Pawel M. Wocjan\\[1ex]Department of Electrical Engineering and Computer Science\\ University of Central Florida\\Orlando}
\email{wocjan@eecs.ucf.edu}

\begin{document}

\maketitle
%%%%%%%%%%%%%%%%%%%%%%%%%%%%%%%%%%%%%%%%%%%%%%%%%%%%%%%%%%%%%%%%%%%%%%%

\begin{slide}{Joint work}
\vspace{-0.2cm}
\begin{itemize}
\item \emph{Quantum Algorithms for One-Dimensional Infrastructures}

\medskip
with Pradeep Sarvepalli (University of British Columbia) 
 
\medskip
preprint arXiv:1106.6347

\bigskip
\item \emph{Quantum Algorithm for Computing the Period Lattice of an Infrastructure}
 
\medskip
with Felix Fontein (University of Z\"urich)
 
\medskip
preprint to appear soon on the arXiv

\bigskip
\item P.W. gratefully acknowledges the support from the NSF grant CCF-0726771 and the NSF CAREER Award CCF-0746600
\end{itemize}
\end{slide}

%%%%%%%%%%%%%%%%%%%%%%%%%%%%%%%%%%%%%%%%%%%%%%%%%%%%%%%%%%%%%%%%%%%%%%%%%

\begin{slide}{Motivation and overview}
\vspace{-0.3cm}
\begin{itemize}
\item infrastructures are group-like algebraic objects

\medskip
\item they appear in arithmetic geometry in the study of computational problems related to
algebraic number fields and function fields with finite constant fields

\medskip
\item arithmetic geometry provides a unified treatment of global fields 

\medskip
\item
the most prominent computational tasks are computing 

\begin{itemize}
\item the period lattice of the infrastructure and

\medskip
\item generalized discrete logarithms
\end{itemize}

\medskip
\item
both problems are not known to have efficient classical algorithms
\end{itemize}

\end{slide}

%%%%%%%%%%%%%%%%%%%%%%%%%%%%%%%%%%%%%%%%%%%%%%%%%%%%%%%%%%%%%%%%%%%%%%%%%

\begin{slide}{QA for problems in algebraic geometry}
\begin{itemize}
\item there is an elegant and efficient quantum algorithm for infrastructures arising from function fields (of curves over finite fields)

\medskip
\item such infrastructures are discrete and embed into finite abelian groups (divisor class groups of degree zero)

\medskip
\item there are efficient (classical) methods to do arithmetics in these groups

\medskip
\item $\Rightarrow$ we can simply use the quantum algorithm for the abelian HSP to compute the period lattice of discrete infrastructures

\end{itemize}

\end{slide}

%%%%%%%%%%%%%%%%%%%%%%%%%%%%%%%%%%%%%%%%%%%%%%%%%%%%%%%%%%%%%%%%%%%%%%%%%

\begin{slide}{QA for problems in algebraic geometry}
\vspace{-0.2cm}
\begin{itemize}
\item the quantum algorithm for the abelian HSP can also be used to efficiently

\medskip
\begin{itemize}
\item compute generalized DLOGs in discrete infrastructures

\medskip
\item compute the whole divisor group class and the ideal class group

\medskip
\item solve the principal ideal problem

\medskip
\item compute the Zeta function

\end{itemize}

\medskip
\item $\Rightarrow$ we focus our attention on non-discrete infrastructures

\medskip
\item such infrastructures arise from number fields

\end{itemize}

\end{slide}

%%%%%%%%%%%%%%%%%%%%%%%%%%%%%%%%%%%%%%%%%%%%%%%%%%%%%%%%%%%%%%%%%%%%%%%%%

\begin{slide}{One-dimensional pseudo-periodic states}

\begin{itemize}
\item let $\Lambda=R \mathbb{Z}$ with $R\in\mathbb{R}_{>0}$

\medskip
\item a quantum state in $\C^{qN}$ of the form
\[
|\psi\> \propto \sum_{\lambda\in\Lambda\cap [0,qN)} | v + N \lambda + \xi_\lambda \> \quad\mbox{with $\xi_\lambda\in (-1,1)$}
\]

is called a pseudo-periodic with period $R$ (with period lattice $\Lambda$)

\medskip
\item $v\in\{0,1,\ldots,\lfloor N R \rfloor\}$ is random

\medskip
\item there are approximately $\tfrac{q}{\mathrm{det}(\Lambda)} = \tfrac{q}{R}$ elements in $|\psi\>$

\medskip
\item task: find a rational number $\hat{R}$ such that $|\hat{R}-R|\le \gamma$
\end{itemize}

\end{slide}

%%%%%%%%%%%%%%%%%%%%%%%%%%%%%%%%%%%%%%%%%%%%%%%%%%%%%%%%%%%%%%%%%%%%%%%%%

\begin{slide}{QA for estimating the period}

\begin{itemize}

\item INPUT:

\begin{itemize}
\item two positive integer $q$ and $N$ with $q > 2 R$ and $N > \tfrac{2}{R}$

\medskip
\item two pseudo-periodic states in $\mathbb{C}^{qN}$ with period $R$
\end{itemize}

\medskip
\item embed each state into a register $\mathbb{C}^{2qN}$

\medskip
\item apply the QFT of size $2qN$ to each state

\medskip
\item measure both registers

\medskip
\item denote the outcomes by $w_1$ and $w_2$

\end{itemize}
\end{slide}

%%%%%%%%%%%%%%%%%%%%%%%%%%%%%%%%%%%%%%%%%%%%%%%%%%%%%%%%%%%%%%%%%%%%%%%%%

\begin{slide}{Sampling approximations of $\Lambda^*$}

\begin{itemize}
\item let $\Lambda^* = \tfrac{1}{R} \mathbb Z$ be the dual lattice of $\Lambda= R \mathbb{Z}$

\medskip
\item let $\mathcal{Y}$ be the set of $w\in\{0,1,\ldots,2qN\}$ satisfying

\medskip
\begin{itemize}
\item $w_i = [2q \lambda^*]$ and for some $\lambda^*\in\Lambda^*$ and 

\medskip
\item $0\le w_i \le 2q (\kappa N)$ with $\kappa\in (0,1)$
\end{itemize}

\medskip
\item for all $w\in\mathcal{Y}$, we have
\[
\Pr(w) \ge \frac{1}{2qN} \cos^2\big( \pi (\tfrac{1}{4} + \tfrac{1}{4qN} + 2\kappa) \big) \frac{q}{R}
\]

\medskip
\item the cardinality of $\mathcal{Y}$ is approximately $\tfrac{\kappa N}{\mathrm{det}(\Lambda^*)} = \kappa N R$

\end{itemize}

\end{slide}

%%%%%%%%%%%%%%%%%%%%%%%%%%%%%%%%%%%%%%%%%%%%%%%%%%%%%%%%%%%%%%%%%%%%%%%%%

\begin{slide}{Approximate generating set for $\Lambda^*$}
\begin{itemize}
\item Let $a$ and $b$ be two numbers chosen uniformly at random from $\{1,2,\ldots,B\}$ for $B\ge 2$

\[
\Pr\big( \gcd(a,b) = 1 \big) \ge 1 - \sum_{p \, \mathrm{prime}} \frac{1}{p^2} > \frac{1}{2}
\]

\medskip
\item for half of the pairs $(w_1,w_2)\in\mathcal{Y}\times\mathcal{Y}$, the corresponding $\lambda_i^*$ generate $\Lambda^*$

\medskip
\item $\Rightarrow$ the pair $(\tfrac{w_1}{2q},\tfrac{w_2}{2q})$ is an approximate generating set of $\Lambda^*$ with probability at least
\[
\frac{\kappa^2}{8} \cos^4\big( \pi (\tfrac{1}{4} + \tfrac{1}{4qN} + 2\kappa) \big) \ge 10^{-5} \quad \mbox{for } \kappa=\tfrac{1}{16}
\]
\end{itemize}

\end{slide}

%%%%%%%%%%%%%%%%%%%%%%%%%%%%%%%%%%%%%%%%%%%%%%%%%%%%%%%%%%%%%%%%%%%%%%%%%

\begin{slide}{Approximate basis for $\Lambda$}

\begin{itemize}

\item we can recover an approximate basis of $\Lambda^*$ from the approximate generating set formed by $\tfrac{w_1}{2q}$ and $\tfrac{w_1}{2q}$  

\medskip
\item we can recover an approximate basis of $\Lambda^*$, i.e., $\hat{R}$ by taking the inverse of the approximate basis of $\Lambda^*$

\end{itemize}

\end{slide}

%%%%%%%%%%%%%%%%%%%%%%%%%%%%%%%%%%%%%%%%%%%%%%%%%%%%%%%%%%%%%%%%%%%%%%%%%

\begin{slide}{Success probability of the QA}

\begin{itemize}
\item our algorithm yields the approximation $\hat{R}$ with probability at least $10^{-5}$ 

\medskip
\item Hallgren's analysis yields the significantly worse bound $\Omega(1/\log^4(M))$ where $M$ is an upper bound on $RN$ and 
works only if $\log(M) \ge 256$

\medskip
\[
\frac{1}{256^4} = \frac{1}{4\,294\,967\,296} < \frac{1}{4 \cdot 10^9}
\]

\bigskip
$\Rightarrow$ the guaranteed expected running time would be at least $50$ days if it took $1$ ms to execute the quantum algorithm once

\end{itemize}
\end{slide}

%%%%%%%%%%%%%%%%%%%%%%%%%%%%%%%%%%%%%%%%%%%%%%%%%%%%%%%%%%%%%%%%%%%%%%%%%

\begin{slide}{One-dimensional infrastructures}
\begin{itemize}

\medskip
\item the problem of computing the period lattice of a one-dimensional infrastructure
\[ 
\mbox{{\Large $\Downarrow$}}
\]
the problem of estimating the period of a one-dimensional periodic quantum state

\medskip
\item one-dimensional infrastructures arise from number fields of unit rank one

\medskip
\item 
quadratic number fields $\mathbb{Q}(\sqrt{d})$ with $d$ positive and squarefree are only one type of number fields of unit rank one
\end{itemize}

\end{slide}

%%%%%%%%%%%%%%%%%%%%%%%%%%%%%%%%%%%%%%%%%%%%%%%%%%%%%%%%%%%%%%%%%%%%%%%%%

\begin{slide}{Improvements}
\vspace{-0.4cm}
\begin{itemize}
\item quadratic number fields yield one-dimensional IS\\ $\Rightarrow$ 
our algorithms can be used to solve Pell's equation and the principal ideal problem

\medskip
\item in this sense, they generalize Hallgren's QA

\medskip
\item our improvements and extensions:

\medskip
\begin{itemize}
\item give a conceptually simpler presentation 

\medskip
\item obtain an improved technique of analyzing Fourier sampling of one-dimensional pseudo-periodic states

\medskip
\item find algorithms with lower complexity and constant success probability
\end{itemize}
\end{itemize}
\end{slide}

%%%%%%%%%%%%%%%%%%%%%%%%%%%%%%%%%%%%%%%%%%%%%%%%%%%%%%%%%%%%%%%%%%%%%%%%%

\begin{slide}{One-dimensional infrastructure}

\begin{itemize}
\item a one-dimensional infrastructure of circumference $R$ is a pair $(X,d)$ where

\begin{itemize}
\item $X=\{x_0,x_1,\ldots,x_{m-1}\}$ is a finite set and

\medskip
\item $d \, : \, X \rightarrow \R/R\Z$ is injective
\end{itemize}

\medskip
\item the elements of $X$ are ordered such that 
\[
0=d(x_0)<d(x_1)<\ldots<d(x_{m-1})<R
\]
we use representatives in $[0,R)$ to order the $d(x_i)$

\medskip
\item imagine that we place the elements around a circle of circumference $R$ starting at the top most point and moving clockwise
\end{itemize}

\end{slide}

%%%%%%%%%%%%%%%%%%%%%%%%%%%%%%%%%%%%%%%%%%%%%%%%%%%%%%%%%%%%%%%%%%%%%%%%%

\begin{slide}{Baby step}

\begin{itemize}
\vspace{-0.25cm}
\item the baby step $\mathrm{bs} : X \rightarrow X$ is defined as
\[
\mathrm{bs}(x_i) = \left\{
\begin{array}{ll}
x_{i+1} & 0 \le i < m-1 \\
x_0     & i=m-1
\end{array}
\right.
\]
\medskip
it gives you the next element on the circle

\medskip
\item the relative distance function $\Delta_{\mathrm{bs}} : X \rightarrow [0,R)$ is such that
\[
d(\mathrm{bs}(x)) = d(x) + \Delta_{\mathrm{bs}}(x)
\]
holds for all $x\in X$

\medskip
it tells you how far you have to travel along the circle to get to the next element
\end{itemize}

\end{slide}

%%%%%%%%%%%%%%%%%%%%%%%%%%%%%%%%%%%%%%%%%%%%%%%%%%%%%%%%%%%%%%%%%%%%%%%%%

\begin{slide}{Giant step}
\vspace{-0.25cm}
\begin{itemize}
\item for $x,y\in X$ consider the set
\[
S_{x,y} = \{ s \in [0,R) \, | \, d(x) + d(y) + s \in d(X) \}
\]

\medskip
\item the giant step $\mathrm{gs} : X \times X \rightarrow X$ is defined as
\[
\mathrm{gs}(x,y)=z 
\]

such that
\[
d(z) = d(x) + d(y) + \min S_{x,y}
\]

\medskip
\item the relative distance function $\Delta_{\mathrm{gs}} \, : \, X \times X \rightarrow [0,R)$ is defined as
\[
\Delta_{\mathrm{gs}}(x,y) = \min S_{x,y}
\]
\end{itemize}

\end{slide}

%%%%%%%%%%%%%%%%%%%%%%%%%%%%%%%%%%%%%%%%%%%%%%%%%%%%%%%%%%%%%%%%%%%%%%%%%

\begin{slide}{Computational problems}
\vspace{-0.25cm}
\begin{itemize}
\item estimate the circumference $R$

\medskip
\item given $x\in X$, estimate $d(x)$

\bigskip
\item these problems are computationally hard for classical computers

\medskip
\item the first generalizes the problem of computing the order of a (black-box) cyclic group

\medskip
\item the second generalizes the problem of determining discrete logarithms in a (black-box) cyclic group
\end{itemize}
\end{slide}

%%%%%%%%%%%%%%%%%%%%%%%%%%%%%%%%%%%%%%%%%%%%%%%%%%%%%%%%%%%%%%%%%%%%%%%%%

\begin{slide}{Finite cyclic groups}
\vspace{-0.25cm}
\begin{itemize}
\item suppose $G=\<g\>$ is a finite cyclic group of order $R$

\medskip
\item $\Rightarrow$ we can form an infrastructure out of $G$ as follows:

\begin{itemize}
\item $X=G$

\medskip
\item $d(h) = \log_g(h)$ 

\medskip
\item the baby step corresponds to multiplication by $g$

\medskip
\item the giant step corresponds to the multiplication of two elements in $G$

\medskip
\item $\Delta_{\mathrm{bs}}$ always takes on the value $1$

\medskip 
\item $\Delta_{\mathrm{gs}}$ always takes on the value $0$
\end{itemize}
\end{itemize}
\end{slide}

%%%%%%%%%%%%%%%%%%%%%%%%%%%%%%%%%%%%%%%%%%%%%%%%%%%%%%%%%%%%%%%%%%%%%%%%%

\begin{slide}{Group arithmetic with $f$-reps}

\begin{itemize}
\item a pair $(x,t)\in X\times \R$ is called an $f$-representation if 
\[
0 \le t < \Delta_{\mathrm{bs}}(x)
\]

\medskip
\item denote the set of all $f$-representations by $\mathrm{Rep}(\mathcal{I})$

\medskip
\item $\mathrm{Rep}(\mathcal{I})$ can be endowed with a group structure such that 
\[
\mathrm{Rep}(\mathcal{I}) \cong \R/R\Z
\]

\medskip
\item set
\[
\Lambda = R \mathbb{Z}
\]
\end{itemize}
\end{slide}

%%%%%%%%%%%%%%%%%%%%%%%%%%%%%%%%%%%%%%%%%%%%%%%%%%%%%%%%%%%%%%%%%%%%%%%%%

\begin{slide}{HSP over $\R$}

\begin{itemize}
\item we have a surjective group homomorphism

\[
f : \R \rightarrow \mathrm{Rep}(\mathcal{I}) \quad \mbox{with}\quad
\ker f = \Lambda
\]

\medskip
\item $\Rightarrow$ intuition: determine the subgroup $\Lambda$ of $\R$ 
hidden by $f$

\medskip
\item but there are two fundamental problems:

\medskip
\begin{itemize}
\item we can only evaluate at a finite number of points

\medskip
\item we cannot compute $f$ exactly due to finite precision since $\Delta_{\mathrm{bs}}$ and $\Delta_{\mathrm{gs}}$
can only be approximated
\end{itemize}

\item $\Rightarrow$ we have to work with a finite number of carefully chosen evaluation points that avoid ``forbidden intervals''
\end{itemize}

\end{slide}

%%%%%%%%%%%%%%%%%%%%%%%%%%%%%%%%%%%%%%%%%%%%%%%%%%%%%%%%%%%%%%%%%%%%%%%%%

\begin{slide}{Computable version}

\vspace{-0.3cm}
\begin{itemize}
\item Let $h : \mathcal{V} \rightarrow X \times \mathcal{V}$ be defined as
\[
h(v) = (x,\lfloor N t \rfloor)
\]

\begin{itemize}
\item $\mathcal{V}=\{0,\ldots,qN-1\}$

\medskip
\item $s$ is some random offset 

\medskip
\item $f(s+\tfrac{v}{N}) = (x,t)$
\end{itemize}

\medskip
\item we can compute $h(v)=(x,k)$ correctly for all $v\in \mathcal{V}$ with high probability

\medskip
\item for most $(x,k)\in h(\mathcal{V})$, the preimage has the special form
\[
h^{-1}(x,k) = \{ v + N \lambda + \xi_\lambda \, : \, \lambda\in\Lambda \} \subseteq \mathcal{V}, \quad \xi_j\in (-1,1) 
\]
\end{itemize}

\end{slide}

%%%%%%%%%%%%%%%%%%%%%%%%%%%%%%%%%%%%%%%%%%%%%%%%%%%%%%%%%%%%%%%%%%%%%%%%%

\begin{slide}{Reduction to pseudo-periodic states}

\begin{itemize}

\item evaluate $h$ in superposition over $\mathcal{V}$ and measure

\medskip
\item this process prepares a pseudo-periodic state in $\C^{qN}$

\[
|\psi\> \propto \sum_{\lambda\in R\mathbb{Z} \cap [0,qN)} 
| v + N \lambda + \xi_\lambda \> \quad\mbox{with $\xi_\lambda\in (-1,1)$}
\]

\medskip
\item $\Rightarrow$ we can apply our quantum algorithm to obtain the estimate $\hat{R}$
 
\end{itemize}

\end{slide}


%%%%%%%%%%%%%%%%%%%%%%%%%%%%%%%%%%%%%%%%%%%%%%%%%%%%%%%%%%%%%%%%%%%%%%%%%

\begin{slide}{Higher-dim. pseudo-periodic states}
\vspace{-0.25cm}
\begin{itemize}
\item let $\Lambda$ be a full-rank lattice in $\R^n$

\medskip
\item a quantum state in ${(\mathbb{C}^{qN})}^{\otimes n}$ of the form

\[
\sum_{\lambda\in\Lambda\cap [0,qN)^n} | v + N \lambda + \xi_\lambda\>, \quad \xi_\lambda\in (-1,1)^n
\]

is called pseudo-periodic state with period lattice $\Lambda$

\medskip
\item
Task: find an approximate basis $\lambda'_1,\ldots,\lambda'_n$ of $\Lambda$

\medskip
\item 
$\lambda'_1,\ldots,\lambda'_n$ is called a $\gamma$-approximate basis iff

\begin{itemize}
\medskip
\item $\Lambda=\mathbb{Z}\lambda_1 + \ldots + \mathbb{Z} \lambda_n$ 

\medskip
\item $\|\lambda_i - \lambda'_i \|_2 \le \gamma$
\end{itemize}
\end{itemize}

\end{slide}

%%%%%%%%%%%%%%%%%%%%%%%%%%%%%%%%%%%%%%%%%%%%%%%%%%%%%%%%%%%%%%%%%%%%%%%%%

\begin{slide}{Higher-dimensional infrastructures}

\begin{itemize}
	\item an $n$-dimensional infrastructure~$\mathcal{I}$ is a finite set of distinguished
  points on an $n$-dimensional torus $\mathbb{R}^n / \Lambda$
  
  \medskip
  \item $\Lambda$ is a lattice of full rank in $\R^n$ 
  
  \medskip
  \item to each distinguished point $x$, we assign a region on the torus, so
  that every point on the torus lies in exactly one such region
  
  \medskip
  \item $\Rightarrow$ every point~$y$ in the region of $x$ can be represented by the difference $t := y - x$ together
  with $x$, i.e.\ as the pair $(x, t)$ 
  
  \medskip
  \item these tuples $(x, t)$ are called the $f$-representations
  of the infrastructure
\end{itemize}

\end{slide}

%%%%%%%%%%%%%%%%%%%%%%%%%%%%%%%%%%%%%%%%%%%%%%%%%%%%%%%%%%%%%%%%%%%%%%%%%

\begin{slide}{Group arithmetic with $f$-reps.}

\begin{itemize}
\item denote the set of all $f$-representations by $\mathrm{Rep}(\mathcal{I})$

\bigskip
\item $\mathrm{Rep}(\mathcal{I})$ can be endowed with a group structure such that 
\[
\mathrm{Rep}(\mathcal{I}) \cong \R^n / \Lambda
\]
\end{itemize}

\end{slide}

%%%%%%%%%%%%%%%%%%%%%%%%%%%%%%%%%%%%%%%%%%%%%%%%%%%%%%%%%%%%%%%%%%%%%%%%%

\begin{slide}{HSP over $\R^n$}

\begin{itemize}
\item exactly as in the one-dimensional case, we have a surjective group homomorphism

\[
f : \R^n \rightarrow \mathrm{Rep}(\mathcal{I}) \quad \mbox{with}\quad
\ker f = \Lambda
\]

\medskip
\item $\Rightarrow$ intuition: determine the subgroup $\Lambda$ of $\R^n$ 
hidden by $f$

\medskip
\item but there are two fundamental problems:

\medskip
\begin{itemize}
\item we can only evaluate at a finite number of points

\medskip
\item we cannot compute $f$ exactly because we cannot can only approximate $f$-representations 
\end{itemize}

\item $\Rightarrow$ we have to work with a finite number of carefully chosen evaluation points that avoid ``forbidden regions''
\end{itemize}

\end{slide}

%%%%%%%%%%%%%%%%%%%%%%%%%%%%%%%%%%%%%%%%%%%%%%%%%%%%%%%%%%%%%%%%%%%%%%%%%%%%

\begin{slide}{Forbidden regions}

\end{slide}

%%%%%%%%%%%%%%%%%%%%%%%%%%%%%%%%%%%%%%%%%%%%%%%%%%%%%%%%%%%%%%%%%%%%%%%%%%%%

\begin{slide}{Reduction to pseudo-periodic states}

\begin{itemize}
\item evaluate $h$ in superposition over $\mathcal{V}=\{0,\ldots,qN-1\}^n$ and measure 

\medskip
\item with high probability, we obtain the pseudo-periodic state

\[
\sum_{\lambda\in\Lambda\cap [0,qN)^n} | v + N \lambda + \xi_\lambda\>, \quad \xi_\lambda\in (-1,1)^n
\]

\end{itemize}

\end{slide}

%%%%%%%%%%%%%%%%%%%%%%%%%%%%%%%%%%%%%%%%%%%%%%%%%%%%%%%%%%%%%%%%%%%%%%%%%

\begin{slide}{QA for estimating the period lattice}

\begin{itemize}

\item INPUT:

\begin{itemize}
\item two positive integers $q$ and $N$ that are sufficiently large

\medskip
\item $2n+1$ pseudo-periodic states in ${(\mathbb{C}^{qN})}^{\otimes n}$ with period lattice $\Lambda$
\end{itemize}

\medskip
\item embed each state into a register ${(\mathbb{C}^{2nqN})}^{\otimes n}$

\medskip
\item apply the multidimensional QFT to each state

\medskip
\item measure 

\medskip
\item denote the outcomes by $w_1,\ldots,w_{2n+1}$

\end{itemize}
\end{slide}

%%%%%%%%%%%%%%%%%%%%%%%%%%%%%%%%%%%%%%%%%%%%%%%%%%%%%%%%%%%%%%%%%%%%%%%%%

\begin{slide}{Sampling approximations of $\Lambda^*$}

\begin{itemize}
\item let $\Lambda^*$ be the dual lattice of $\Lambda$

\medskip
\item let $\mathcal{Y}$ be the set of $w\in\{0,1,\ldots,2nqN-1\}^n$ satisfying

\medskip
\begin{itemize}
\item $w_i = [2n q \lambda^*]$ and for some $\lambda^*\in\Lambda^*$ and 

\medskip
\item $\|w_i \|_\infty \le 2n q (\kappa N)$ with $\kappa\in (0,1)$
\end{itemize}

\medskip
\item for all $w\in\mathcal{Y}$, we have
\[
\Pr(w) \ge \frac{1}{(2qnN)^n} \cos^2\big( \pi (\tfrac{1}{4} + \tfrac{1}{4qN} + 2\kappa) \big) \frac{q^n}{\mathrm{det}(\Lambda)}
\]

\medskip
\item the cardinality of $\mathcal{Y}$ is approx. $\tfrac{(\kappa N)^n}{\mathrm{det}(\Lambda^*)} = (\kappa N)^n \mathrm{det}(\Lambda)$

\end{itemize}

\end{slide}

%%%%%%%%%%%%%%%%%%%%%%%%%%%%%%%%%%%%%%%%%%%%%%%%%%%%%%%%%%%%%%%%%%%%%%%%%

\begin{slide}{Approximate generating set for $\Lambda^*$}
\begin{itemize}
\item we sample $n$ vectors of $\Lambda^* \cap [0,b)^n$ 

\medskip
\item with prob. at least $\tfrac{1}{4}$, they generate a full-rank sublattice $\Lambda^*_0$

\medskip
\item $\Lambda^*/\Lambda^*_0$ is a finite group generated by $n$ elements

\medskip
\item we sample $n+1$ vectors of $\Lambda^* \cap [0,b_0)^n$ with $b_0 > b$

\medskip
\item this induces an almost uniform distribution of $\Lambda^*/\Lambda^*_0$

\medskip
\item with prob. at least $\frac{1}{4}-\hat{\zeta}>0.184$ the new and old samples generate $\Lambda^*$ 

\end{itemize}

\end{slide}

%%%%%%%%%%%%%%%%%%%%%%%%%%%%%%%%%%%%%%%%%%%%%%%%%%%%%%%%%%%%%%%%%%%%%%%%%

\begin{slide}{Approximate basis for $\Lambda$}

\begin{itemize}

\item we can recover an approximate basis of $\Lambda^*$ from the approximate generating set formed by 

\[
	\frac{w_1}{2nq},\ldots,\frac{w_{2n+1}}{2nq}
\]  

\medskip
\item we can recover an approximate basis of $\Lambda^*$from the approximate basis of $\Lambda^*$

\medskip
\item the algorithm succeeds with probability at least
\[
	\frac{1}{4} \Big( \frac{1}{4}-\hat{\zeta} \Big) \left( \frac{c \kappa^n}{(2n)^n} \right)^{2n+1}
\]

where $c=\cos^2\big( \pi ( \tfrac{1}{4} + \tfrac{1}{4qN} + 2\kappa) \big)$ and $\kappa = \tfrac{1}{16n}$
\end{itemize}

\end{slide}

%%%%%%%%%%%%%%%%%%%%%%%%%%%%%%%%%%%%%%%%%%%%%%%%%%%%%%%%%%%%%%%%%%%%%%%

\begin{slide}{Improvements}

\begin{itemize}

\item we present a rigorous and complete proof that the running time is polynomial in the determinant of the period lattice $\Lambda$

\medskip
\item we give an explicit lower bound on the success probability 

\medskip
{\sc Hallgren} and {\sc Schmidt and Vollmer} did not give any lower bound

\medskip
{\sc Schmidt} presented a lower bound in his dissertation which is worse than our bound  

\medskip
\item message: it is very important to improve the lower bound ! 

\medskip
we are sure that this can be done !

\end{itemize}

\end{slide}

\begin{slide}{Thank you!!!}

\end{slide}

\end{document} 